\chapter{複素引数実数値関数}
    \section{解析的な複素引数実数値関数は定数関数である}
        \label{解析的な複素引数実数値関数は定数関数である}
        \begin{shadebox}
            $G\subseteq\complexNumbers$を開領域とし、$f:G\to\realNumbers$は$G$上で解析的であるとする。
            このとき、$f$は定数関数である。
        \end{shadebox}
        \noindent
        (証明1)
        \begin{proof}
            Cauchy-Riemannの関係式から直ちに従う。
        \end{proof}
        \noindent
        (証明2)
        \newcommand{\Rh}{\prescript{\text{R}}{}{h}}
        \newcommand{\Ih}{\prescript{\text{I}}{}{h}}
        \begin{proof}
            \quad\par
            $x_0\in G$を任意にとる。
            絶対値が十分小さい$\Rh,\Ih$を$x_0+\Rh, x_0+i\Ih \in G$となるように任意にとる。
            $f$が$x_0$に於いて解析的であるから$f'(x_0)$が存在し、次式が成り立つ。
            \begin{align*}
                \frac{f(x_0+\Rh) - f(x_0)}{\Rh} &= f'(x_0) + \frac{o(|\Rh|)}{\Rh} \tag{1}\\
                \frac{f(x_0+i\Ih) - f(x_0)}{i\Ih} &= f'(x_0) + \frac{o(|\Ih|)}{i\Ih} \tag{2}
            \end{align*}
            式(1),(2)の左辺はそれぞれ実数, 純虚数であるので、$\Rh\to 0,\Ih\to 0$の極限を考えると、$f'(x_0) = 0$でなくてはならない
            (背理法で容易に示せる)。
            $x_0$は$G$内の任意の点であったから、$f$は$G$上で定数関数である。
        \end{proof}
    \section{1次増分}
        \newcommand{\RG}{\prescript{\text{R}}{}{G}}
        \newcommand{\Rf}{\prescript{\text{R}}{}{f}}
        \begin{shadebox}
            $G\subseteq\complexNumbers$を開領域とし、$\RG\coloneqq\setParens{(x,y)\in\realNumbers^2}{x+iy\in G}$とする。
            $\Rf:\RG\to\realNumbers$は$\RG$上で$\mathrm{C}^1$級であるとし、$f:z\in G\mapsto\Rf\parens{\Re{z},\Im{z}}\in\realNumbers$とする。
            $z\in G$を任意にとり、絶対値が十分小さい$h\in G$を$z+h\in G$となるように任意にとる。
            $f(z+h)-f(z)$の$h$に関する1次の項は、Cauchy-Riemannの関係式から導かれる微分係数を形式的に計算して得られる$f$の形式的な微分係数$f'(z)$と$h$の積の実部$\Re{f'(z)h}$である。
        \end{shadebox}
        \begin{proof}
            \quad\par
            $x=\Re{z},\;y=\Im{z},\;\Rh=\Re{h},\Ih=\Im{h}$とすると次式が成り立つ。
            \begin{align*}
                &\phantom{=} f(z+h) - f(z) \\
                &= \Rf(x+\Rh,y+\Ih) - \Rf(x,y) = \partDeriv{\Rf}{x}{}(x,y)\Rh + \partDeriv{\Rf}{y}{}(x,y)\Ih + o(\sqrt{\Rh^2+\Ih^2}) \\
                &= \frac{1}{2}\left[\partDeriv{\Rf}{x}{}(x,y)\Rh - i\partDeriv{\Rf}{y}{}(x,y)(\Ih i)\right] + \frac{1}{2}\left[\partDeriv{\Rf}{x}{}(x,y)\Rh + i\partDeriv{\Rf}{y}{}(x,y)(-\Ih i)\right] + o(|h|) \\
                &= \frac{1}{2}\underbrace{\left[\partDeriv{\Rf}{x}{}(x,y) - i\partDeriv{\Rf}{y}{}(x,y)\right]}_{\text{(1)}}h + \frac{1}{2}\underbrace{\left[\partDeriv{\Rf}{x}{}(x,y) + i\partDeriv{\Rf}{y}{}(x,y)\right]}_{\text{(2)}}\overline{h} + o(|h|)
            \end{align*}
            式(1)は複素引数実数値関数に対して複素微分係数の実関数表示を形式的に適用して計算したものと一致する。
            今、これを形式的に$f'(z)$と書くことにする。
            なお、\ref{解析的な複素引数実数値関数は定数関数である}で述べている通り、複素引数実数値関数で解析的なものは常に定数関数であるため、形式的に導入した$f'(z)$は本来の複素微分の意味とは似て非なるものである。
            \[ f(z+h) - f(z) = \frac{1}{2}f'(z)h + \frac{1}{2}\overline{f'(z)h} + o(|h|) = \Re{f'(z)h} + o(|h|)\]
        \end{proof}
    \let\Rh\undefined
    \let\Ih\undefined
    \let\Rf\undefined
    \section{複素引数凸関数}
        \subsection{複素引数凸関数の例}
            \label{複素引数凸関数の例}
            \subsubsection{\texorpdfstring{$\|\sum_{i=1}^N X_i\bm{a}_i + \bm{b}\|_2^2\;(\bm{a}_i, \bm{b} \in \complexNumbers^n,\;X_i \in \complexNumbers^{m\times n})$}{行列X1,X2,...,XNの線形結合×ベクトル+定数のノルム}は\texorpdfstring{$X_1,\dotsc,X_N$}{X1,X2,...,XN}に関して凸である。}
                \begin{proof}
                    \quad\par
                    $T_X \coloneqq (X_1,\dotsc,X_N),\;T_Y \coloneqq (Y_1,\dotsc,Y_N),\;\lambda \in [0,1]$とし、$f(T_X) \coloneqq \|\sum_{i=1}^N X_i\bm{a}_i + \bm{b}\|_2^2$とする。
                    \begin{align*}
                        f((1-\lambda)T_X + \lambda T_Y) &= \left\|\sum_{i=1}^N ((1-\lambda)X_i + \lambda Y_i)\bm{a}_i + \bm{b}\right\|_2^2 \\
                        &= \left\|(1-\lambda)\left(\sum_{i=1}^N X_i\bm{a}_i + \bm{b}\right) + \lambda\left(\sum_{i=1}^N Y_i\bm{a}_i + \bm{b}\right)\right\|_2^2 \\
                        &= (1-\lambda)^2\left\|\sum_{i=1}^N X_i\bm{a}_i + \bm{b}\right\|_2^2 + \lambda^2\left\|\sum_{i=1}^N Y_i\bm{a}_i + \bm{b}\right\|_2^2 \\
                        &\phantom{=} + 2\lambda(1-\lambda)\Re{\left(\sum_{i=1}^N X_i\bm{a}_i + \bm{b}\right)^\mathrm{H}\left(\sum_{i=1}^N Y_i\bm{a}_i + \bm{b}\right)}
                    \end{align*}
                    よって
                    \begin{align*}
                        &\phantom{=} (1-\lambda)f(T_X) + \lambda f(T_Y) - f((1-\lambda)T_X + \lambda T_Y) \\
                        &= \lambda(1-\lambda)\left[\left\|\sum_{i=1}^N X_i\bm{a}_i + \bm{b}\right\|_2^2 + \left\|\sum_{i=1}^N Y_i\bm{a}_i + \bm{b}\right\|_2^2 - 2\Re{\left(\sum_{i=1}^N X_i\bm{a}_i + \bm{b}\right)^\mathrm{H}\left(\sum_{i=1}^N Y_i\bm{a}_i + \bm{b}\right)}\right] \\
                        &= \lambda(1-\lambda)\left\|\left(\sum_{i=1}^N X_i\bm{a}_i + \bm{b}\right) - \left(\sum_{i=1}^N Y_i\bm{a}_i + \bm{b}\right)\right\|_2^2 \geq 0
                    \end{align*}
                \end{proof}
            \subsubsection{\texorpdfstring{$\norm{A\bm{x}+\bm{b}}_2^2$}{||Ax+b||\string^2}が狭義凸\texorpdfstring{$\iff A^*A\succ O$}{⇔ A*Aが正定}}
                \begin{shadebox}
                    $A \in \complexNumbers^{m\times n}, \bm{x} \in \complexNumbers^n, \bm{b} \in \complexNumbers^m$とする。
                    $f(\bm{x}) \coloneqq \norm{A\bm{x}+\bm{b}}_2^2$が$\bm{x}$に関して狭義凸となる必要十分条件は$A^*A\succ O$である。
                \end{shadebox}
                $\lambda \in (0,1),\; \bm{x},\bm{y} \in \complexNumbers^m, \bm{x}\neq\bm{y}$とする。
                \begin{proof}
                    まず次式が成り立つ。
                    \begin{align*}
                        f((1-\lambda)\bm{x} + \lambda\bm{y}) &\coloneqq \norm{(1-\lambda)(A\bm{x}+\bm{b}) + \lambda(A\bm{y}+\bm{b})}_2^2 \\
                        &= (1-\lambda)^2 f(\bm{x}) + \lambda^2 f(\bm{y}) + 2\lambda(1-\lambda)\Re{(A\bm{x}+\bm{b})^*(A\bm{y}+\bm{b})}
                    \end{align*}
                    これより
                    \begin{align*}
                        &\phantom{=} (1-\lambda)f(\bm{x}) + \lambda f(\bm{y}) - f((1-\lambda)\bm{x} + \lambda\bm{y}) \\
                        &= \lambda(1-\lambda)\left[f(\bm{x}) + f(\bm{y}) - 2\Re{(A\bm{x}+\bm{b})^*(A\bm{y}+\bm{b})}\right] \\
                        &= \lambda(1-\lambda)\norm{(A\bm{x}+\bm{b}) - (A\bm{y}+\bm{b})}_2^2 = \lambda(1-\lambda)\norm{A(\bm{x}-\bm{y})}_2^2 \\
                        &= \lambda(1-\lambda)(\bm{x}-\bm{y})^*A^*A(\bm{x}-\bm{y})
                    \end{align*}
                    $\lambda \in (0,1),\;\bm{x}\neq\bm{y}$に対して上式が正となる必要十分条件は$A^*A\succ O$である。
                \end{proof}
